\documentclass[10pt,a4paper]{amsart}
\usepackage[margin=19mm]{geometry}
\usepackage{amsmath,amssymb,mathtools}
\usepackage[scr=rsfs,cal=euler]{mathalfa}
\usepackage{xfrac}
\usepackage[unicode,hidelinks,bookmarks=false]{hyperref}
\newcommand{\ie}{i.e.\ }
\newcommand{\cf}{cf.\ }
\newcommand{\resp}{respectively\ }
\newcommand{\Subsection}{\subsection}
\providecommand{\nobreakdash}{\nobreak\hbox{-}}
\setlength{\emergencystretch}{2em}
\multlinegap0pt
\usepackage{multirow}
\usepackage[shortlabels]{enumitem}
\usepackage{booktabs}
\usepackage{xcolor}
\usepackage{tikz}
\usepackage[normalem]{ulem}
\usepackage{amssymb}
\usepackage{braket}
\usepackage{caption}
\newtheorem{lemma}{Lemma}
\title{\textbf{Conclusive Identification Via Noisy Classical Channel: Superactivation and Quantum Advantage}}
\author{Anushko Chattopadhyay, Ambuj, Rakesh Das, Smritikana Patra, Chitrak Roychowdhury, Manik Banik, Amit Mukherjee}
\date{Source: arXiv:2604.00089v1 (CC BY 4.0)}
\begin{document}
\maketitle

\noindent\emph{Source and licence.} \textbf{Conclusive Identification Via Noisy Classical Channel: Superactivation and Quantum Advantage}. Version arXiv:2604.00089v1 (CC BY 4.0), distributed by arXiv under the Creative Commons Attribution 4.0 International licence, http://creativecommons.org/licenses/by/4.0/, as recorded in the arXiv licence metadata for this exact version and shown on the abstract page https://arxiv.org/abs/2604.00089v1. Full licence notice: \texttt{licenses/arxiv-2604.00089-CC-BY-4.0.txt}.\par\smallskip

\noindent\textit{Editorial scope.} Counted unit: the article's lemma lemma1 (source0.tex lines 474--476). The article's complete original proof of it is delivered with it (source0.tex lines 478--523, one \texttt{proof} environment of 46 lines). No mathematics is written, repaired or re-derived: the statement and the proof are the article's own lines, in the article's own notation, and no retained line is substituted or re-flowed.  Retained uncounted as the source-local context the proof needs: lines 367--420.  Nothing was omitted from any retained span, so the package carries no omission record.  No retained line contains a citation, so the wrapper carries no bibliography and no bibliography entry.  No retained line contains a figure environment, an image-inclusion command or drawing code, so no image is delivered and the package declares no asset dependency.  Editorial formatting only, in addition: an AMS \texttt{amsart} preamble that restates the article's own declarations and macros used by the retained lines, namely the environments \texttt{lemma} on separate counters, together with the standard packages those lines need.  The retained environments print in this document as Lemma 1 in order of appearance, whereas the article prints the same items with the numbering of its own full document.  The complete native source of the article as distributed in the arXiv e-print for this exact version is bundled unchanged as \texttt{sources/197563/source0.tex}.
\noindent Notably, a support graph corresponds to a unique confusability graph, but the same confusability graph can arise from different support graphs, and hence from different channels. For instance, consider the SNFC channels $N_i : X \to X$ for $i \in \{1, \cdots, 6\}$ with $X := \{1,\cdots, 5\}$, whose adjacency matrices (denoted $\mathbb{S}_{N_i}=\mathbb{S}_i$) are given by
\begin{align}
{\fontsize{8}{9}\selectfont
\begin{aligned}
\mathbb{S}_1 = 
\begin{pmatrix}
        1 & 1 & 0 & 0 & 1 \\
        1 & 1 & 1 & 0 & 0 \\
        0 & 1 & 1 & 1 & 0 \\
        0 & 0 & 1 & 1 & 1 \\
        1 & 0 & 0 & 1 & 1
\end{pmatrix};~~
\mathbb{S}_2 = 
\begin{pmatrix}
        1 & 1 & 1 & 0 & 1 \\
        1 & 1 & 1 & 0 & 0 \\
        1 & 1 & 1 & 1 & 0 \\
        0 & 0 & 1 & 1 & 1 \\
        1 & 0 & 0 & 1 & 1
\end{pmatrix};~~
\mathbb{S}_3= 
\begin{pmatrix}
        1 & 1 & 1 & 1 & 1 \\
        1 & 1 & 1 & 0 & 0 \\
        1 & 1 & 1 & 1 & 0 \\
        1 & 0 & 1 & 1 & 1 \\
        1 & 0 & 0 & 1 & 1
\end{pmatrix};\\
\mathbb{S}_4 = 
\begin{pmatrix}
        1 & 1 & 1 & 1 & 1 \\
        1 & 1 & 1 & 0 & 1 \\
        1 & 1 & 1 & 1 & 0 \\
        1 & 0 & 1 & 1 & 1 \\
        1 & 1 & 0 & 1 & 1
\end{pmatrix};~~
\mathbb{S}_5 = 
\begin{pmatrix}
        1 & 1 & 1 & 1 & 1 \\
        1 & 1 & 1 & 1 & 0 \\
        1 & 1 & 1 & 1 & 0 \\
        1 & 1 & 1 & 1 & 1 \\
        1 & 0 & 0 & 1 & 1
\end{pmatrix};~~
\mathbb{S}_6= 
\begin{pmatrix}
        1 & 1 & 1 & 1 & 1 \\
        1 & 1 & 1 & 1 & 1 \\
        1 & 1 & 1 & 1 & 0 \\
        1 & 1 & 1 & 1 & 1 \\
        1 & 1 & 0 & 1 & 1
\end{pmatrix}.
\end{aligned}}\label{penta}
\end{align}

\begin{lemma}\label{lemma1}
For an SNFC channel $N : X \to X$ with $|X| = 5$ whose support graph is $\mathtt{S}_1$ of Eq.~(\ref{penta}), the receiver can conclusively identify all $5$ inputs when assisted with a perfect classical channel $\mathrm{id}^c_3$, i.e. $\mathrm{ci}_\circ(N \otimes \mathrm{id}^c_3) = 5$.
\end{lemma}

\begin{proof}
Let the channel $N$ has input (and output) $X=\{1,\cdots,5\}$, the inputs (outputs) of $\mathrm{id}^c_3$ channel are $\{red (R), green (G), blue (B)\}$. The combined channel $N \otimes \mathrm{id}^c_3$ has input and alphabet $X \times \{R,G,B\}$. Alice partitions $X = \{1,2,3,4,5\}$ into three groups according to the map $\mathcal{P}: X \to \{R,G,B\}$:
\begin{equation*}
\mathcal{P}(2)= \mathcal{P}(4)= B,\quad \mathcal{P}(1) = R, \quad \mathcal{P}(3) = \mathcal{P}(5) = G.
\end{equation*}
Upon receiving the pair $(x, \mathcal{P})\in X\times\{R,G,B\}$, Bob knows which group the input belongs to, and uses the output $x$ of the noisy channel to identify the input conclusively within that group. We verify each group separately in Table~\ref{tab1}.

\begin{table}[t]
\centering
\begin{tabular}{c|c|c|c}
\hline
~Input of $N$~ & ~Output of $N$~ &~ Message via $\mathrm{id}^c_3$~ & ~Identified?~ \\
\hline\hline
\multirow{3}{*}{$1$} 
& $5$ & R & $\checkmark$ \\
& $1$ & R & $\checkmark$ \\
& $2$ & R & $\checkmark$ \\
\hline
\multirow{3}{*}{$2$} 
& $1$ & B & $\checkmark$ \\
& $2$ & B & $\checkmark$ \\
& $3$ & B & $\times$ \\
\hline
\multirow{3}{*}{$3$} 
& $2$ & G & $\checkmark$ \\
& $3$ & G & $\checkmark$ \\
& $4$ & G & $\times$ \\
\hline
\multirow{3}{*}{$4$} 
& $3$ & B & $\times$ \\
& $4$ & B & $\checkmark$ \\
& $5$ & B & $\checkmark$ \\
\hline
\multirow{3}{*}{$5$} 
& $4$ & G & $\times$ \\
& $5$ & G & $\checkmark$ \\
& $1$ & G & $\checkmark$ \\
\hline
\end{tabular}
\caption{The conclusive identification scheme for the SNFC channel $N$ (having support graph $\mathtt{S}_{N}=\mathtt{S}_1$) when assisted by a noiseless classical channel $\mathrm{id}^c_3$. }
\label{tab1}
\end{table}
\noindent We have $\mathrm{ci}_\circ(N) = 0$ and $\mathrm{ci}_\circ(\mathrm{id}^c_3) = 3$, whereas the combined channel $N \otimes \mathrm{id}^c_3$ allows Bob to conclusively identify all $5$ inputs, giving $\mathrm{ci}_\circ(N \otimes \mathrm{id}^c_3) = 5$, 
%Therefore, $\mathrm{ci}_\circ(N \otimes \mathrm{id}^c_3) = 5$, 
which establishes the superadditivity of $\mathrm{ci}_\circ$. This completes the proof. 
\end{proof}

\end{document}
